\section*{Capítulo \ref{ch:b1:electronic-effects} --- Efeitos eletrônicos e intermediários reativos}

\begin{solution}{exo:b1:electronic-effects:1}
2-Metilbutano, \ce{CH3-CH(CH3)-CH2-CH3}: os três \ce{CH3} são primários,
o \ce{CH2} secundário, o CH terciário. 2-Bromo-2-metilpropano: os três
\ce{CH3} são primários, o carbono central terciário. Dos quatro brometos
\ce{C4H9Br} (1-bromobutano, 2-bromobutano, 1-bromo-2-metilpropano,
2-bromo-2-metilpropano), o último dá o \omterm{def:b1:electronic-effects:intermediates}{carbocátion} terciário, o mais
estável.
\end{solution}

\begin{solution}{exo:b1:electronic-effects:2}
$10^{4.76 - 2.87} = 10^{1.89} = 78$; $10^{4.76 - 0.52} = 10^{4.24} =
\num{1.7e4}$.
\end{solution}

\begin{solution}{exo:b1:electronic-effects:3}
Etanoato: a C=O e a \ce{C-O^-} trocadas entre os dois oxigênios.
4-Nitrofenóxido: o \omterm{def:b1:lewis-resonance-vsepr:lewis}{par não ligante} do oxigênio forma uma C=O, as
ligações duplas do anel se deslocam, o carbono que leva \ce{NO2} forma
uma ligação C=N, e um par N=O vai para o oxigênio: uma estrutura com C=O
em uma extremidade do anel, C=N na outra e os dois oxigênios do nitro
negativos (uma estrutura quinoide). Cátion alila:
\ce{CH2=CH-CH2+} $\leftrightarrow$ \ce{+CH2-CH=CH2}.
\end{solution}

\begin{solution}{exo:b1:electronic-effects:4}
Segunda seta: da ligação O--H da água para o seu oxigênio; produtos
\ce{H2O} (vinda de \ce{OH-}, agora neutra) e \ce{OH-} (o oxigênio da
antiga água fica com o par, carga $-1$). Para a amônia: uma seta do \omterm{def:b1:lewis-resonance-vsepr:lewis}{par
não ligante} de N para um H de \ce{H3O+}, e uma dessa ligação O--H para o
oxigênio: N passa a $+1$, O passa a 0.
\end{solution}

\begin{solution}{exo:b1:electronic-effects:5}
Etanol (15.9) $<$ água (14.0) $<$ fenol (9.99, deslocalização no anel)
$<$ 4-nitrofenol (7.16, \ce{NO2} recebe a carga) $<$ ácido etanoico (4.76,
dois oxigênios equivalentes) $<$ ácido tricloroetanoico (0.51, três
cloros $-I$).
\end{solution}

\begin{solution}{exo:b1:electronic-effects:6}
Anilina (4.60) $<$ piridina (5.23) $<$ amônia (9.25) $<$ metilamina
(10.66). Na anilina, o \omterm{def:b1:lewis-resonance-vsepr:lewis}{par não ligante} é compartilhado com o anel
(estruturas com C=N$^+$ e a carga negativa em \textit{orto} e
\textit{para}); protonar N custaria essa deslocalização.
\end{solution}

\begin{solution}{exo:b1:electronic-effects:7}
No 3-nitrofenóxido, as \omterm{def:b1:lewis-resonance-vsepr:resonance}{estruturas de ressonância} põem a carga nos
carbonos \textit{orto} e \textit{para} em relação ao oxigênio, nunca no
carbono que leva \ce{NO2}: nenhuma ajuda mesomérica, daí ser mais fraco
que o isômero 4. O grupo nitro ainda atrai elétrons pelas ligações
$\sigma$ ($-I$), daí ser mais forte que o fenol.
\end{solution}

\begin{solution}{exo:b1:electronic-effects:8}
\ce{CH3+} $<$ \ce{CH3CH2+} $<$ \ce{CH2=CH-CH2+} $<$ \ce{C6H5CH2+}
$\approx$ \ce{(CH3)3C+} $<$ \ce{CH3OCH2+}: os grupos alquila cedem
elétrons; a alila e a benzila deslocalizam a carga (duas e quatro
estruturas); o \omterm{def:b1:lewis-resonance-vsepr:lewis}{par não ligante} do oxigênio dá uma estrutura com octeto
completo em toda parte. A ordem no meio da lista depende do meio.
\end{solution}

\begin{solution}{exo:b1:electronic-effects:9}
\ce{OH-} (\omterm{def:b1:electronic-effects:nucleophile}{nucleófilo}) ataca o carbono de \ce{CH3Br} (\omterm{def:b1:electronic-effects:nucleophile}{eletrófilo}), com o
par C--Br saindo no Br. \ce{H2O} cede um par a \ce{H+} (aqui, ácido e
base são os nomes usados). \ce{CN-} ataca o carbono carbonílico do
etanal, com o par $\pi$ da C=O indo para o oxigênio. \ce{NH3} cede seu
\omterm{def:b1:lewis-resonance-vsepr:lewis}{par não ligante} ao orbital vazio do boro em \ce{BF3}.
\end{solution}

\begin{solution}{exo:b1:electronic-effects:10}
Não: $K_a = 10^{-0.51} = 0.31$ é muito maior que $C$; o ácido está quase
totalmente dissociado. $x^2/(0.010 - x) = 0.31$ dá $x = \qty{9.7e-3}{mol/L}$,
$\mathrm{pH} = 2.01$, \omterm{def:b1:predominant-reaction:dissociation}{grau de dissociação} de \qty{97}{\percent}.
\end{solution}

\begin{solution}{exo:b1:electronic-effects:11}
Os grupos alquila empurram elétrons ($+I$) para o oxigênio do alcóxido e
tornam sua carga menos estável, tanto mais quanto mais grupos houver. Os
alcóxidos maiores também são menos bem solvatados pela água, que
estabiliza melhor um oxigênio carregado pequeno do que um enterrado
entre grupos alquila. Os dois efeitos aumentam o
$\mathrm{p}K_a$.
\end{solution}

\begin{solution}{exo:b1:electronic-effects:12}
Em uma amida, o \omterm{def:b1:lewis-resonance-vsepr:lewis}{par não ligante} do nitrogênio é compartilhado com a C=O
(estrutura \ce{N+=C-O^-}); tomá-lo para um próton ou para um \omterm{def:b1:electronic-effects:nucleophile}{eletrófilo}
destruiria essa estabilização. Em \omterm{def:b1:predominant-reaction:strong-acid}{ácido forte}, a amida é protonada no
oxigênio, que porta a carga parcial negativa.
\end{solution}

\begin{solution}{pb:b1:electronic-effects:1}
\textbf{1.} $10^{1.89} = 78$; $10^{1.52} = 33$; $10^{0.84} = 6.9$.
\textbf{2.} Não: cada cloro adicional ajuda menos. A carga do carboxilato
já está bem espalhada, e as medidas ficam menos precisas à medida que o
ácido se aproxima da dissociação total.
\textbf{3.} Cada Cl atrai densidade eletrônica através das ligações C--C
e C--O e estabiliza o carboxilato ($-I$).
\textbf{4.} O \omterm{def:b1:electronic-effects:inductive}{efeito indutivo} se apaga em duas ou três ligações.
\textbf{5.} Os dois ácidos parecem igualmente fortes (0.52 e 0.51). Na
água, ambos estão quase totalmente dissociados nas concentrações usuais
(\cref{ch:b1:predominant-reaction}, \omterm{def:b1:predominant-reaction:levelling}{efeito nivelador}): um
$\mathrm{p}K_a$ perto de 0 é o limite do que a água consegue distinguir,
e os dois valores têm incertezas de alguns décimos.
\textbf{6.} Ácido etanoico: forma ácida (\qty{98}{\percent});
cloroetanoico: as duas, com o ânion ligeiramente predominante
($10^{3 - 2.87} = 1.3$); os três outros: ânion.
\textbf{7.} \ce{CH3-C(=O)-O^-} $\leftrightarrow$ \ce{CH3-C(-O^-)=O}.
\textbf{8.} Duas ligações C--O iguais, entre uma dupla e uma simples.
\textbf{9.} A carga fica no único oxigênio, sem ligação $\pi$ para
compartilhá-la.
\textbf{10.} $10^{15.9 - 4.76} = 10^{11.1}$, cerca de $10^{11}$.
\textbf{11.} A carga do fenóxido se espalha por três carbonos do anel,
menos eletronegativos que o oxigênio: um ganho, menor que em um
carboxilato: $\mathrm{p}K_a$ 9.99, entre 4.76 e 15.9.
\textbf{12.} $K = 10^{6.37 - \mathrm{p}K_a}$: ácido etanoico $10^{1.61} = 41$,
favorecida (libera-se dióxido de carbono com excesso de
hidrogenocarbonato); fenol $10^{-3.62}$ e etanol $10^{-9.5}$: nenhuma
reação.
\textbf{13.} 4-Nitrofenol (7.16) $>$ 2-nitrofenol (7.23) $>$ 3-nitrofenol
(8.36) $>$ fenol (9.99).
\textbf{14.} A estrutura quinoide do \cref{exo:b1:electronic-effects:3},
com a carga em um oxigênio do grupo nitro.
\textbf{15.} A carga só chega aos carbonos \textit{orto} e \textit{para}
em relação ao oxigênio; o carbono do nitro está em \textit{meta}.
\textbf{16.} O efeito $-I$ de \ce{NO2}.
\textbf{17.} No 2-nitrofenol, o OH forma uma \omterm{def:b1:intermolecular-forces-solvents:hbond}{ligação de hidrogênio} com um
oxigênio do grupo nitro vizinho, o que estabiliza a forma ácida e torna
o próton um pouco mais difícil de retirar.
\textbf{18.} Frações de ânion $1/(1 + 10^{\mathrm{p}K_a - 7.5})$: 4-nitro
\qty{69}{\percent}, 2-nitro \qty{65}{\percent}, 3-nitro
\qty{12}{\percent}. O ácido e o ânion do 4-nitrofenol absorvem de modo
diferente, o ânion no visível: sua cor aparece em torno de seu
$\mathrm{p}K_a$.
\textbf{19.} $x^2/(0.010 - x) = 10^{-0.52} = 0.30$: $x = \qty{9.7e-3}{mol/L}$,
$\mathrm{pH} = 2.01$: quase um \omterm{def:b1:predominant-reaction:strong-acid}{ácido forte}.
\textbf{20.} $\frac12(4.76 + 2.00) = 3.38$.
\textbf{21.} \ce{CF3COOH + CH3COO- -> CF3COO- + CH3COOH}, $K = 10^{4.76 -
0.52} = 10^{4.24}$: \omterm{def:b1:extent-q-and-k:quantitative}{quantitativa}.
\textbf{22.} $K_a(\ce{CF3COOH})/K_a(\ce{CH3COOH}) =$ \textbf{$10^{4.24} =
\num{1.7e4}$}, cerca de dez mil.
\end{solution}
